\chapter{Validation and Testing}
\label{chap:validation}

\section{Why Validation Matters}

A key objective of \actsim is to provide transparent and reproducible
actuarial simulations. Without independent validation, simulation
results are difficult to defend in audit, regulatory review, or
peer review. This chapter describes the validation strategies that
\actsim users should employ to verify that simulations behave as
expected.

The validation strategies described here cover frequency, severity,
aggregate losses, claim development, triangle reconciliation, and
reproducibility. Together, they form a baseline that should be
sufficient for open-source review and CAS-style scrutiny.

\section{Frequency Distribution Validation}

For frequency validation, simulated claim counts should be compared
against the theoretical expected count.

For example, if claim frequency follows
\[
    N \sim \mathrm{Poisson}(\lambda),
\]
then $E[N] = \lambda$ and $\mathrm{Var}(N) = \lambda$. The simulated
mean and variance of the claim counts should converge to $\lambda$ as
the number of simulations $M$ increases. The Monte Carlo standard error
on the simulated mean is approximately $\sqrt{\lambda / M}$.

If a validation example uses $\lambda = 50$ with $M = 100{,}000$
simulations, the simulated average claim count should be close to $50$,
with a standard error of roughly $\sqrt{50/100000} \approx 0.022$. A
simulated mean within a few standard errors of $50$ is consistent with
correct implementation.

For negative binomial frequencies, the analogous checks compare the
simulated mean and variance against the theoretical $E[N]$ and
$\mathrm{Var}(N)$, where the variance exceeds the mean.

\section{Severity Distribution Validation}

For severity validation, the simulated samples should match the moments
and quantiles of the target distribution. For lognormal severity with
parameters $\mu$ and $\sigma$,
\[
    E[X] = \exp\!\left( \mu + \frac{\sigma^{2}}{2} \right),
    \qquad
    \mathrm{Var}(X) = \big( \exp(\sigma^{2}) - 1 \big) \cdot \exp(2\mu + \sigma^{2}).
\]

The empirical mean, variance, and selected percentiles of the simulated
severities should be close to these theoretical values. Larger sample
sizes are required for accurate matching of upper percentiles.

A typical severity validation script:

\begin{lstlisting}[style=pythonstyle]
import numpy as np
from actstats import actuarial as act

mu, sigma = 10, 0.5
n = 1_000_000

samples = act.lognormal(mu, sigma).rvs(size=n)

theoretical_mean = np.exp(mu + sigma**2 / 2)
print("Theoretical mean:", theoretical_mean)
print("Simulated mean:",   samples.mean())
print("Simulated 95%:",    np.percentile(samples, 95))
\end{lstlisting}

\section{Aggregate Loss Validation}

The compound expectation
\[
    E[S] = E[N] \cdot E[X]
\]
provides a simple consistency check on aggregate loss simulations. The
variance under independence is
\[
    \mathrm{Var}(S) = E[N] \cdot \mathrm{Var}(X)
                    + \mathrm{Var}(N) \cdot (E[X])^{2}.
\]
Both quantities should be matched (within sampling error) by simulation
output for a sufficient number of simulations.

When dependency is introduced via correlation or copulas, the simulated
mean should remain close to $E[N] \cdot E[X]$ if the marginal
distributions of $N$ and $X$ are unchanged. The variance, however,
should change in a direction consistent with the imposed dependency.

\section{Claim Development Validation}

For each simulated claim with ultimate loss $U$ and a deterministic
LDF schedule $\{\mathrm{LDF}_a\}$, the expected reported loss at age
$a$ is
\[
    E[L_a] = \frac{U}{\mathrm{CDF}_a},
\]
where $\mathrm{CDF}_a$ is the cumulative product of LDFs from age $a$
to ultimate.

When stochastic LDFs are used, the same expectation holds approximately,
because the multiplicative noise has mean $1$. The simulated reported
losses at each age should average to the deterministic pattern, with
the spread controlled by the volatility parameter.

A simple validation check:

\begin{lstlisting}[style=pythonstyle]
import pandas as pd

# After running simulate_claim_development:
expected_pattern = (
    claim_sim.claim_development
    .groupby("development_month")["incurred_loss"]
    .mean()
)
print(expected_pattern)
\end{lstlisting}

The mean reported loss should grow monotonically with development age
and should approach the simulated ultimate at the latest age.

\section{Triangle Reconciliation}

Once the simulated claim development data have been aggregated into a
loss triangle and analysed using \texttt{chainladder}, the projected
ultimates can be reconciled against the underlying simulated ultimates.

\begin{lstlisting}[style=pythonstyle]
ultimate_simulated = (
    claim_sim.claim_data
    .groupby(claim_sim.claim_data["incurred_date"].dt.year)["ultimate_loss"]
    .sum()
)

ultimate_chainladder = model.ultimate_

print("Simulated ultimate:",   ultimate_simulated)
print("Chain-ladder ultimate:", ultimate_chainladder)
\end{lstlisting}

For a well-specified simulation with a sufficient number of claims, the
two should agree within sampling error. Persistent differences indicate
either a misspecification of the LDF schedule or a sample size that is
too small for the chain-ladder method to converge.

\section{Reproducibility Using Random Seeds}

\actsim runs are fully deterministic given a random seed. The same
input (policy table, distributional specification, seed) produces the
same output across runs. This guarantee is achieved by passing the seed
through to \texttt{numpy}'s random number generator at construction.

To verify reproducibility:

\begin{lstlisting}[style=pythonstyle]
from actsim import StochasticSimulator

def run(seed):
    sim = StochasticSimulator(
        "poisson", (10,), "lognormal", (10, 0.5),
        num_sim=10000, keep_all=False, seed=seed,
    )
    sim.gen_agg_simulations()
    return sim.results.sum()

assert run(1234) == run(1234)
\end{lstlisting}

When sharing simulation results, the seed used to produce them should
be recorded alongside the output.

\section{Statistical Tests}

In addition to mean-and-variance checks, the user may apply formal
statistical tests:

\begin{itemize}[leftmargin=*]
    \item Kolmogorov--Smirnov test on simulated severities against the
          target CDF.
    \item Chi-square goodness-of-fit on simulated frequencies binned
          into integer counts.
    \item Anderson--Darling test for tail behaviour.
\end{itemize}

These tests are most informative when validating a particular
distributional implementation. They are less useful for assessing the
overall realism of a model relative to historical data, where domain
judgement is essential.

\section{Known Limitations}

The current implementation has several limitations that users should be
aware of:

\begin{itemize}[leftmargin=*]
    \item The Kolmogorov--Smirnov statistic is computed via
          \texttt{scipy.stats.kstest} on the data and the fitted CDF.
          For large sample sizes, this can take longer to compute than
          AIC or BIC.
    \item Chi-square statistics in \texttt{DistributionFitter} are
          computed by binning observed and expected frequencies; the
          choice of bin count affects the resulting statistic.
    \item The current copula implementation supports a fixed set of
          families (Gaussian, Frank, Gumbel, Clayton). Other families
          can be added but require a code change.
    \item Multivariate dependency simulation currently supports
          marginal distributions but does not retain claim-level data
          beyond the marginal totals. This limits some downstream
          validations.
\end{itemize}

These limitations are documented to allow reviewers to focus their
scrutiny on the parts of the package that may behave unexpectedly. The
maintainer welcomes contributions that address them.
